\section{The three-factor FX model}\label{sec: threeFactor}
\subsection{The Three Factor model with Local Volatility }
The model is fully described by the following SDEs:
\begin{align*}
 dS_t   = &\ (r^d_t - r^f_t)S_t dt + \sigma(t,S_t)S_tdW^S_t \\
 dr_t^d  = &\  \mu^d_t dt + \sigma_t^d dB_t^d \numberthis \label{eq: threeFactorModel}\\
 dr_t^f  = & \ \big(\mu^f_t - \rho_{Sf}\sigma(t,S_t)\sigma_t^f \big)dt + \sigma_t^f dB_t^f
\end{align*}
where $S_t$ is the spot, $r_t^d$ and $r_t^f$ the domestic and foreign interest rates respectively.\\
The volatility of the spot FX rate is the local volatility function $\sigma(t,S_t)$.
At this stage the dynamics for the two interest rates are left unspecified to let the framework as general as possible. The stochastic vector $(W^S_t, B_t^d, B_t^f)$ is a Brownian Motion under the domestic risk-neutral measure $\mathbb{Q}^d$ with correlation matrix
\[R=\left( \begin{array}{ccc}
1 & \rho_{Sd} & \rho_{Sf} \\
\rho_{Sd} & 1 & \rho_{df} \\
\rho_{Sf} & \rho_{df} & 1 \end{array} \right)\]
The term $-\rho_{Sf}\sigma(t,S_t)\sigma^f_t$ in the drift of the foreign rate is the quanto adjustment: under the foreign risk-neutral measure $\mathbb{Q}^f$ the foreign rate has drift $\mu^f_t$.

\paragraph{Measures.} The domestic and foreign discount factors are $ D_{0t}^i=\exp\big(-\int_0^t r_s^ids\big)$ for $i\in\{d,f\}$. The foreign risk-neutral measure is defined by (section 2.8 and p.256 in \cite{clark})
\begin{equation}\label{eq: RNforeign}
\left.\frac{d\mathbb{Q}^f}{d\mathbb{Q}^d}\right|_{\mathcal{F}_t}=\frac{S_tD^d_{0t}}{S_0D^f_{0t}}
=\exp\bigg( \int_0^t\sigma(s,S_s) dW^S_s - \frac{1}{2}\int_0^t\sigma(s,S_s)^2 ds\bigg),
\end{equation}
so that, denoting with a tilde the Brownian motions under $\mathbb{Q}^f$,
\begin{align*}
dW^S_t = &\ d\widetilde{W}^{S}_t + \sigma(t,S_t)dt \\
dB^d_t = &\ d\widetilde{B}^{d}_t + \rho_{Sd}\sigma(t,S_t)dt  \numberthis \label{measureChange: domesticToForeign}\\
dB^f_t = &\ d\widetilde{B}^{f}_t + \rho_{Sf}\sigma(t,S_t)dt
\end{align*}
Note that the spot driver is shifted by the full $\sigma(t,S_t)$, and that the Girsanov kernel is a function of the state. We write $\mathbb{Q}^{d^T}$ and $\mathbb{Q}^{f^T}$ for the domestic and foreign $T$-forward measures, with numeraires $P^d_{tT}$ and (in foreign currency) $P^f_{tT}$.

\subsection{The Calibration Condition}
In the proposition below, we give a necessary and sufficient condition for a three factor SIR-LVm to be calibrated to a given market smile. $C(t,K)$ denotes the fair market value of a European call option with strike $K$ and maturity $t$.
From \cite{deelstraRayee}, the Dupire volatility is of the form:
\begin{equation}\label{eq: threeFactorDupireVolatility}
\sigma_{Dup}^2(t,K) =
\frac{  \partial_t C(t,K) + (r_d^0 - r_f^0)K\partial_K C(t,K)+ r_f^0 C(t,K)}
{ \frac{1}{2}K^2\partial^2_K C(t,K) }
\end{equation}
where $r_i^0 = r_i^0(t)= \mathbb{E}^{i^t} [r^i_t] = -\partial_t\ln P^{i}_{0t}$ is the market instantaneous forward rate, $\mathbb{E}^{i^t}$ being the expectation under the $t$-forward domestic and foreign measures for $i = d, i=f$ respectively.

\begin{prop}[Calibration of 3-Factor SIR-LVm to market smiles]\label{prop: calibration}\ \\
Assume the model exists with a $C^{1,2}$ call surface and sufficient
integrability for the discounted Tanaka formula below. Assume a smooth market
surface with positive strike density, matching initial and boundary conditions,
and uniqueness for the linear forward equation with the model's fixed local
variance coefficient. Under these hypotheses the model
(\ref{eq: threeFactorModel}) is exactly calibrated to the market smile if and only if
\begin{align*}
\sigma(t,K)^2 = \sigma_{Dup}(t,K)^2 & - \frac{ \mathbb{E}^d[ D_{0t}^d(r_t^d-r^0_d)\mathbb{I}_{S_t>K} ]}{ \frac{1}{2}K\partial^2_KC(t,K) }\\
 &+ S_0\frac{ \mathbb{E}^f[ D_{0t}^f(r_t^f-r^0_f)\mathbb{I}_{S_t>K} ]}{ \frac{1}{2}K^2\partial^2_KC(t,K) }\numberthis \label{eq: threeFactorCalibrationCondition}
\end{align*}
or, equivalently, under the domestic measure only,
\begin{equation}\label{eq: singleMeasure}
\sigma(t,K)^2 = \sigma_{Dup}(t,K)^2 - \frac{ \mathbb{E}^d\big[ D_{0t}^d\big(K(r_t^d-r^0_d)-S_t(r_t^f-r^0_f)\big)\mathbb{I}_{S_t>K} \big]}{ \frac{1}{2}K^2\partial^2_KC(t,K) }
\end{equation}
\end{prop}

\begin{proof}\ \\
We follow the proof of Proposition 11.1 in \cite{guyonLabordere} and show that the result is an easy generalisation of (\ref{eq:calibration_condition_spot}). \\
Define $\mathcal{P}_t\equiv D^d_{0t}(S_t-K)^+$ to be the discounted vanilla call payoff with maturity $t$ and strike $K$ under the domestic measure. By applying Ito-Tanaka's formula on such payoff, we get:

\begin{align*}
d\mathcal{P}_t  = & - r_t^d D^d_{0t}(S_t-K)^+dt + D^d_{0t}S_t\mathbb{I}_{S_t>K} \{ (r_t^d-r_t^f)dt + \sigma(t,S_t)dW^S_t \}\\
& + \frac{1}{2} D^d_{0t}\delta_K(S_t)S^2_t\sigma(t,S_t)^2dt \\
= &  - r_t^d D^d_{0t}S_t\mathbb{I}_{S_t>K}dt + r_t^d D^d_{0t}K\mathbb{I}_{S_t>K}dt + r_t^d D^d_{0t}S_t\mathbb{I}_{S_t>K}dt\\
&  -r_t^f D^d_{0t}S_t\mathbb{I}_{S_t>K}dt + D^d_{0t}\mathbb{I}_{S_t>K}\sigma(t,S_t)S_tdW^S_t+ \frac{1}{2} D^d_{0t}\delta_K(S_t)K^2\sigma(t,K)^2dt\\
= & D^d_{0t}\mathbb{I}_{S_t>K}(r_t^d K- r_t^f S_t)dt +  D^d_{0t}\mathbb{I}_{S_t>K}\sigma(t,S_t)S_tdW^S_t 
+ \frac{1}{2} D^d_{0t}\delta_K(S_t)\sigma(t,K)^2K^2 dt
\end{align*}
By taking the expectation $\mathbb{E}^d[\cdot]$ on both sides of the above equation and by assuming that $ M_t =\int_0^tD^d_{0s}\mathbb{I}_{S_s>K}\sigma(s,S_s)S_s dW^S_s$ is a true martingale, we get
$$ \partial_tC_m(t,K)= \mathbb{E}^d\big[ D^d_{0t}\mathbb{I}_{S_t>K}(r_t^d K- r_t^f S_t) \big] + \frac{1}{2}\sigma(t,K)^2K^2 \mathbb{E}^d\big[ D^d_{0t}\delta_K(S_t) \big]$$
where $C_m(t,K)= \mathbb{E}^d[\mathcal{P}_t]$ is the price of the call option in the model.\\
By using that $\partial_K^2C_m(t,K)= \mathbb{E}^d\big[ D^d_{0t}\delta_K(S_t) \big] $, we rewrite
\begin{align*}
\partial_tC_m(t,K) = &\ \mathbb{E}^d\big[ D^d_{0t}\mathbb{I}_{S_t>K}(r_t^d K- r_t^f S_t) \big] + \frac{1}{2}\sigma(t,K)^2K^2\partial_K^2C_m(t,K)\\
= &\ \mathbb{E}^d\big[ D^d_{0t}\mathbb{I}_{S_t>K}(r_t^d - r_d^0 )K \big] + \mathbb{E}^d\big[ D^d_{0t}\mathbb{I}_{S_t>K} r_d^0 K \big] \\
&-\mathbb{E}^d\big[ D^d_{0t}\mathbb{I}_{S_t>K}(r_t^f - r_f^0 )S_t \big] - \mathbb{E}^d\big[ D^d_{0t}\mathbb{I}_{S_t>K} r_f^0 S_t \big]\\
& + \frac{1}{2}\sigma(t,K)^2K^2\partial_K^2C_m(t,K)
\end{align*}

We now rewrite the term in $r^f$ under the foreign measure. From section 2.8 in \cite{clark}, one has that the Radon–Nikodym derivative of $\mathbb{P}^f$ with respect to $\mathbb{P}^d$ is
$$ \frac{d\mathbb{P}^f}{d\mathbb{P}^d}=\exp\bigg( \int_0^t\sigma(s,S_s) dW^S_s - \frac{1}{2}\int_0^t\sigma(s,S_s)^2 ds\bigg)$$
Hence:
\begin{align*}
& \mathbb{E}^d\big[ D^d_{0t}\mathbb{I}_{S_t>K}(r_t^f - r_f^0 )S_t \big] \\
& = \mathbb{E}^d\bigg[ D^d_{0t}\mathbb{I}_{S_t>K}(r_t^f - r_f^0 )S_0\exp\bigg( \int_0^t (r_s^d - r_s^f)ds + \int_0^t\sigma(s,S_s) dW^S_s - \frac{1}{2}\int_0^t\sigma(s,S_s)^2 ds \bigg) \bigg] \\
& =  S_0\mathbb{E}^d\bigg[\mathbb{I}_{S_t>K}(r_t^f - r_f^0 )\exp\bigg( -\int_0^t r_s^fds\bigg)\cdot\exp\bigg( \int_0^t\sigma(s,S_s) dW^S_s - \frac{1}{2}\int_0^t\sigma(s,S_s)^2 ds \bigg)  \bigg] \\
& =  S_0\mathbb{E}^d\bigg[D^f_{0t}\mathbb{I}_{S_t>K}(r_t^f - r_f^0 )\frac{d\mathbb{P}^f}{d\mathbb{P}^d}\bigg]\\
& =  S_0\mathbb{E}^f\big[D^f_{0t}\mathbb{I}_{S_t>K}(r_t^f - r_f^0 )\big]
\end{align*}
Together with the fact that $\partial_K C_m(t,K)= -\mathbb{E}^d\big[ D^d_{0t}\mathbb{I}_{S_t>K} \big] $, one has
\begin{align*}
\partial_tC_m(t,K) = & \ K\mathbb{E}^d\big[ D^d_{0t}\mathbb{I}_{S_t>K}(r_t^d - r_d^0 )\big] -  S_0\mathbb{E}^f\big[D^f_{0t}\mathbb{I}_{S_t>K}(r_t^f - r_f^0 )\big]\\
&- K \partial_K C_m(t,K)r_d^0 - r_f^0\mathbb{E}^d[D^d_{0t}\mathbb{I}_{S_t>K} S_t]+ \frac{1}{2}\sigma(t,K)^2K^2\partial_K^2C_m(t,K)\\
= & \ K\mathbb{E}^d\big[ D^d_{0t}\mathbb{I}_{S_t>K}(r_t^d - r_d^0 )\big] -  S_0\mathbb{E}^f\big[D^f_{0t}\mathbb{I}_{S_t>K}(r_t^f - r_f^0 )\big]\\
& - r_d^0 K \partial_K C_m(t,K) - r_f^0\mathbb{E}^d[D^d_{0t}(S_t-K)^+]- r_f ^0K\mathbb{E}^d[D^d_{0t}\mathbb{I}_{S_t>K}]\\
& + \frac{1}{2}\sigma(t,K)^2K^2\partial_K^2C_m(t,K)\\
= & \ K\mathbb{E}^d\big[ D^d_{0t}\mathbb{I}_{S_t>K}(r_t^d - r_d^0 )\big] -  S_0\mathbb{E}^f\big[D^f_{0t}\mathbb{I}_{S_t>K}(r_t^f - r_f^0 )\big]\\
& - (r_d^0 - r_f^0) K \partial_K C_m(t,K) - r_f^0 C_m(t,K)  \\
& + \frac{1}{2}\sigma(t,K)^2K^2\partial_K^2C_m(t,K)
\end{align*}
So the model is calibrated to the market smile if and only if 
\begin{align*}
\partial_tC(t,K) =  & \ K\mathbb{E}^d\big[ D^d_{0t}\mathbb{I}_{S_t>K}(r_t^d - r_d^0 )\big] -  S_0\mathbb{E}^f\big[D^f_{0t}\mathbb{I}_{S_t>K}(r_t^f - r_f^0 )\big]\\
& - (r_d^0 - r_f^0) K \partial_K C(t,K) - r_f^0 C(t,K)  \\
& + \frac{1}{2}K^2\sigma(t,K)^2\partial_K^2C(t,K)
\end{align*}
The result follows easily from the definition of $\sigma_{Dup}(t,K)$ in (\ref{eq: threeFactorDupireVolatility}). For sufficiency, fix the law of the existing
model and hence its rate-tail forcing term. The asserted identity makes the
market surface solve the same inhomogeneous linear forward equation as $C_m$,
with coefficient $\sigma^2$ and this same forcing. Their difference solves the
homogeneous equation with zero initial and boundary data; the assumed
uniqueness gives $C_m=C$.
\end{proof}

\begin{rem}
The equivalence of (\ref{eq: threeFactorCalibrationCondition}) and (\ref{eq: singleMeasure}) is the identity $D^d_{0t}S_t=S_0D^f_{0t}\,\frac{d\mathbb{Q}^f}{d\mathbb{Q}^d}$ used in the proof. In the case of absence of a foreign rate (i.e. $r^f_t\equiv 0$), we get the usual calibration condition in (\ref{eq:calibration_condition_spot}).
\end{rem}

The form (\ref{eq: singleMeasure}) is the one to use numerically: both rate corrections are expectations under the same measure. Restated under the domestic $T$-forward measure $\mathbb{Q}^{d^T}$ it reads
\begin{equation}\label{eq: threeFactorCalibrationConditionForward}
\sigma(t,K)^2 = \sigma_{Dup}(t,K)^2  - P^d_{0T}\frac{ \mathbb{E}^{\mathbb{Q}^{d^T}}\big[ (P^d_{tT})^{-1}\big(K(r_t^d-r^0_d)-S_t(r_t^f-r^0_f)\big)\mathbb{I}_{S_t>K} \big]}{ \frac{1}{2}K^2\partial^2_KC(t,K) }
\end{equation}

\subsection{The Particle Method}
From Lemma 5.2.2 in \cite{hybridModel}, the model (\ref{eq: threeFactorModel}) under the domestic $T$-forward measure $\mathbb{Q}^{d^T}$ reads:
\begin{align*}
 \frac{df_t}{f_t}   = &\ \sigma(t,S_t)dW^{S^T}_t - \sigma_{P_{tT}}^d dB_t^{d^T} + \sigma_{P_{tT}}^f dB_t^{f^T} \\
 dr_t^d  = &\  (\mu^d_t + \sigma_t^d\sigma_{P_{tT}}^d)dt + \sigma_t^d dB_t^{d^T} \numberthis \label{eq: threeFactorModelForward}\\
 dr_t^f  = & \ \big(\mu^f_t - \rho_{Sf}\sigma(t,S_t)\sigma_t^f +\rho_{df}\sigma_t^f\sigma_{P_{tT}}^d \big)dt + \sigma_t^f dB_t^{f^T}
\end{align*}
where $f_t:=S_t P^f_{tT}/P^d_{tT}$ is the forward and $\sigma_{P_{tT}}^d$, $\sigma_{P_{tT}}^f$ the volatilities of the domestic $P^d_{tT}$ and foreign $P^f_{tT}$ bonds respectively.\\
A \textit{particle} is then described by the vector-valued process $(f_t^i,r_t^{d,(i)},r_t^{f,(i)})$, for $1\leq i \leq N$, where $N$ is the number of simulated particles. Using the representation (\ref{eq: threeFactorCalibrationConditionForward}) for the local volatility, define:
\begin{equation}\label{eq: discreteLocVol3FactorModel}
\sigma_N(t,K)^2  =  \sigma_{Dup}(t,K)^2  - P^d_{0T}\frac{\frac{1}{N}\sum_{i=1}^N (P_{tT}^{d,(i)})^{-1}\big(K(r_t^{d,(i)}-r^0_d)-S^i_t(r_t^{f,(i)}-r^0_f)\big)\mathbb{I}_{S_t^i>K} }{\frac{1}{2}K^2\partial^2_KC(t,K)}
\end{equation}
All particles are simulated under $\mathbb{Q}^{d^T}$: no second system of particles under the foreign forward measure is needed. The estimator (\ref{eq: discreteLocVol3FactorModel}) still has the two drawbacks discussed in Section \ref{subsec: malliavinGeneral}: it depends on $\partial_K^2C$, and it is a ratio of two small quantities far from the money. The Malliavin representation of Section \ref{subsec: malliavin3F} removes both from the rates correction. It does not remove them from the problem: $\sigma_{Dup}$ itself still requires $\partial_K^2C$, and the kernel regression still needs particles near the strike, so whether the wings behave better in a full calibration remains to be tested.

\subsection{Hull White Stochastic Interest Rates}\label{subsec: HW3F}
Assume both the domestic and foreign rates follow Hull-White dynamics as in (\ref{eq: hull_white}), with parameters $(\kappa_i,\sigma^i_r)$, $i\in \{d,f\}$. It is convenient to work with the centred states $x^i_t$ defined by
\begin{equation}\label{eq: centredStates}
r^i_t=r^0_i(t)+\frac{(\sigma^i_r)^2}{2\kappa_i^2}\big(1-e^{-\kappa_it}\big)^2+x^i_t,\qquad x^i_0=0,
\end{equation}
which under $\mathbb{Q}^d$ follow (linear SDEs, see for instance \cite{bruti})
\begin{align*}
dx^d_t&=-\kappa_dx^d_t\,dt+\sigma^d_r\,dB^d_t\\
dx^f_t&=\big(-\kappa_fx^f_t-\rho_{Sf}\sigma^f_r\,\sigma(t,S_t)\big)dt+\sigma^f_r\,dB^f_t
\end{align*}
In particular the centred rates
$$ c^i_t:=r^i_t-r^0_i(t)=x^i_t+\frac{(\sigma^i_r)^2}{2\kappa_i^2}\big(1-e^{-\kappa_it}\big)^2 $$
are known pathwise: the centring constant of the foreign rate is the market forward rate $r^0_f(t)=\mathbb{E}^{f^t}[r^f_t]$, which is \emph{not} $\mathbb{E}^{d^t}[r^f_t]$, and no expectation of $\sigma(s,S_s)$ has to be estimated.

The domestic bond price as a function of the domestic rate, and the foreign bond price as a function of the foreign rate, are given as in (\ref{eq: bondPrice_hullWhite}):
$$ P^i_{tT}= \frac{P_{0T}^{i,mkt}}{P_{0t}^{i,mkt}}\exp\bigg( -A^i(t,T)\,c^i_t-\frac{(\sigma^i_r)^2}{4\kappa_i^3}(e^{-\kappa_i T}-e^{-\kappa_i t})^2(e^{2\kappa_i t}-1) \bigg),\qquad A^i(t,T) = \frac{1-e^{ -\kappa_i(T-t) }}{\kappa_i}, $$
and the volatilities of the bonds are $\sigma_{P_{tT}}^i=-\sigma^i_r A^i(t,T)$. The Brownian motions of $\mathbb{Q}^d$ and $\mathbb{Q}^{d^T}$ are related by
\begin{align*}
dW^{S}_t = &\ dW^{S^T}_t - \rho_{Sd}\sigma^d_r A^d(t,T)dt \\
dB^{d}_t = &\ dB^{d^T}_t - \sigma^d_r A^d(t,T)dt  \numberthis \label{measureChange: domesticForwardToDomestic}\\
dB^{f}_t = &\ dB^{f^T}_t - \rho_{df}\sigma^d_r A^d(t,T)dt
\end{align*}
and the model (\ref{eq: threeFactorModelForward}) now reads, with $S_t=f_tP^d_{tT}/P^f_{tT}$:
\begin{align*}
 \frac{df_t}{f_t}   = &\ \sigma(t,S_t)dW^{S^T}_t + \sigma^d_r A^d(t,T) dB_t^{d^T} - \sigma^f_r A^f(t,T) dB_t^{f^T} \\
 dx_t^d  = &\  \big[-\kappa_dx^d_t - (\sigma^d_r)^2 A^d(t,T)\big]dt + \sigma^d_r dB_t^{d^T} \numberthis \label{eq: threeFactorModelForwardHullWhite}\\
 dx_t^f  = & \ \big[-\kappa_fx^f_t - \rho_{Sf}\sigma^f_r\sigma(t,S_t) - \rho_{df}\sigma^d_r\sigma^f_rA^d(t,T) \big]dt + \sigma^f_r dB_t^{f^T}
\end{align*}
